\section{Results}
\label{sec:results}

\subsection{Saturation audit}

Figure~\ref{fig:audit} summarises the audit. Over \ChiffreAuditN{} Sobol
samples of the v5 box, $\log_{10}k_p$ stays within
$[\ChiffreKpPcinq;\ChiffreKpPqs]$ (5th and 95th percentiles, median
\ChiffreKpMed) and reaches its lower bound in \ChiffreKpSat{} of cases. The
$79~\%$ saturation reported by the earlier campaign would require
$k_p<10^{-4}$, four orders of magnitude below what the projection produces.
It therefore signals a defect \emph{upstream} of the clipping, not a bound set
too tight.

Two components saturate at $100~\%$: $\log_{10}\omega_f$ and $b$. Both are
\emph{constants} of the synthesis ($\omega_f=1/(r_CC)$ and $b=1$ by
relation~\eqref{eq:projection}): the saturation is structural and removes no
degree of freedom. The FO-PIDF vector produced by the synthesis thus has five
effective degrees of freedom, not seven.

\begin{figure}[t]
\centering
\includegraphics[width=0.9\linewidth]{fig_audit_saturation.pdf}
\caption{Saturation audit of the anchoring bounds (figure labels in French).
No varying component saturates; the two observed saturations affect
constants of the synthesis.}
\label{fig:audit}
\end{figure}

\subsection{Exploratory probe: the feasible region exists}

Before the campaign, an exploratory probe (seeds 1 to 6, distinct from the
campaign seeds) searched jointly over the whole of $\theta$, with 1500 calls
per method. The nominally feasible region \emph{exists}: with $b$ free, the
best margin is \ChiffreMargeBLibre{} after replay at $0.2~\mu$s
(\ChiffreNFaisBLibre{} methods out of \ChiffreNLancersLibre{} reach it); with
$b=1$ imposed, it is still \ChiffreMargeBUn{} (\ChiffreNFaisBUn{} out of
\ChiffreNLancersUn), figure~\ref{fig:leverb}. The parameter $b$ is therefore
a margin lever, not a feasibility barrier, contrary to the hypothesis we first
put forward. However, the feasible point found with $b=1$ requires
$\log_{10}k_p\approx-1.66$, outside the image of the v5 box under the
synthesis; the one found with $b$ free requires $b=0$, which v5 cannot set.

\begin{figure}[t]
\centering
\includegraphics[width=0.9\linewidth]{fig_levier_b.pdf}
\caption{Exploratory probe: best nominal margin per method, $b$ free versus
$b=1$. Joint search over the whole of $\theta$, 1500 calls per run.}
\label{fig:leverb}
\end{figure}

\subsection{Confirmatory campaign}

Table~\ref{tab:arms} summarises the three arms. \textbf{None of the 540 cells
is robustly feasible on the design set}, and the Wilson interval per method
and per arm is $[0;\ChiffreWilsonHaut]$.

\begin{table}[t]
\centering
\caption{Confirmatory campaign, 30 seeds $\times$ 6 methods per arm.}
\label{tab:arms}
\begin{tabular}{lccc}
\toprule
 & v5 (brief) & v5b ($b$ free) & v5kp ($k_p$ free) \\
\midrule
cells & \ChiffreNCinq & \ChiffreNCinqB & \ChiffreNCinqKp \\
sterile seeds (stage 1) & \ChiffreSterilesCinq & \ChiffreSterilesCinqB & \ChiffreSterilesCinqKp \\
cells with a nominally feasible point & \ChiffreNomFaisCinq & \ChiffreNomFaisCinqB & \ChiffreNomFaisCinqKp \\
best nominal margin (stage 1) & \ChiffreMargeEtUnCinq & \ChiffreMargeEtUnCinqB & \ChiffreMargeEtUnCinqKp \\
robustly feasible cells & \ChiffreRobFaisCinq & \ChiffreRobFaisCinqB & \ChiffreRobFaisCinqKp \\
best robust margin (stage 2) & \ChiffreMargeEtDeuxCinq & \ChiffreMargeEtDeuxCinqB & \ChiffreMargeEtDeuxCinqKp \\
Friedman, $p$ & \ChiffreFriedmanPCinq & \ChiffreFriedmanPCinqB & \ChiffreFriedmanPCinqKp \\
significant pairs (Holm, of 15) & \ChiffrePairesCinq & \ChiffrePairesCinqB & \ChiffrePairesCinqKp \\
method with best median margin & \ChiffreMeilleureCinq & \ChiffreMeilleureCinqB & \ChiffreMeilleureCinqKp \\
\bottomrule
\end{tabular}
\end{table}

The v5 arm, as specified by the brief, finds \emph{no} nominally feasible
point in 700 calls, for any of the 180 cells. Freeing $b$ (v5b) makes nominal
feasibility reachable in \ChiffreNomFaisCinqB{} cells out of 180, and freeing
$k_p$ (v5kp) in \ChiffreNomFaisCinqKp; in both cases, no cell holds on the
design set.

Friedman tests are significant in every arm. These differences, however,
concern \emph{proximity} to feasibility (all margins are negative), not
feasibility achieved. The ranking of methods changes from one arm to another
(\ChiffreMeilleureCinq{} leads on v5, \ChiffreMeilleureCinqB{} on v5b): no
universal superiority is established (figure~\ref{fig:seeds}).

\begin{figure}[t]
\centering
\includegraphics[width=0.49\linewidth]{fig_graines_v5.pdf}\hfill
\includegraphics[width=0.49\linewidth]{fig_graines_v5b.pdf}
\caption{Robust margin per seed (the statistical unit), arms v5 and v5b.
Symlog scale: linear over $[-1;1]$. Figure labels in French.}
\label{fig:seeds}
\end{figure}

\paragraph{Stage 3 and sealed test.}
Since no candidate is feasible on the design set, stage~3 has nothing to
replay and the sealed test \emph{was not played}. It remains untouched,
seed \ChiffreGraineTest{} included, and keeps its full probative value for
a later study.

\subsection{Cross-validation against the Simulink model}

Six frozen candidates were replayed on the reference Simulink model (nominal
plant), with the same controller, translated into an S-function, and the same
cost and constraint formulas. Agreement criteria were fixed before the run.
The \ChiffreXvalN{} candidates with a controller agree: difference in $J$ of at
most \ChiffreXvalDJmax, difference in each $g_i$ of at most \ChiffreXvalDgmax,
identical feasibility verdict. Both feasible points of the probe remain
feasible under Simulink (margin \ChiffreXvalMargeSl{} for the $b=1$ point).

The feedback-free witness fails the criterion on $g_4$ (difference
\ChiffreXvalFFdg): its error on the 1~V plateau differs by \ChiffreXvalFFde{}
between the two models. Without feedback, nothing corrects the
discontinuous-conduction equilibrium, which depends on the diode and switch
models; in closed loop, the integrator removes this difference. The same cause
explains a $V_\mathrm{min}$ that is on average \ChiffreXvalDVmin{} higher under
Simulink. The floor at $d_\mathrm{min}$ remains below the threshold there
($V_\mathrm{min}=\ChiffreXvalVminSl$): the physical obstruction is confirmed by
the reference model.

\subsection{Two post hoc diagnostics}

The following two diagnostics were added \emph{after} observation. They are
reported, but were used for no selection.

\paragraph{Voltage floor at minimum duty ratio.}
The $g_6$ margins recur identically, to the fourth decimal, across
controllers and arms. At $d=d_\mathrm{min}=0.02$ in open loop, the converter
in discontinuous conduction settles at \ChiffrePlancherNominal{} in the
nominal case, and below $0.8$~V in \ChiffrePlancherN{} of the
\ChiffrePlancherTot{} cases (down to \ChiffrePlancherMin). Any controller that
stays saturated at the floor during the $24\to1$~V descent is trapped there
and violates $g_6$; to satisfy $g_6$, it must leave saturation before the
voltage falls below $0.8$~V.

\paragraph{Limit cycles invisible to the constraints.}
The best controller of the probe, feasible with respect to the six
constraints, sustains at the 24~V plateau a limit cycle at about $7.7$~kHz,
of 624~mV peak-to-peak on $v_o$ averaged over a PWM period
(figure~\ref{fig:traj}). The controller obtained with $b=1$ also oscillates,
more weakly (158~mV). Plateau errors are window means, which cancel the
oscillation, and the $|\Delta d|$ term of $J$ is too weak to penalise it.
With $b=0$, the output also drops from about 24 to 14~V at start-up, a
transient that no constraint detects.

\begin{figure}[t]
\centering
\includegraphics[width=0.85\linewidth]{fig_trajectoire_b_libre.pdf}
\caption{Trajectory of the best controller of the probe ($b$ free). Feasible
with respect to the six constraints, it shows a start-up dip and a limit
cycle at the 24~V plateau. Figure labels in French.}
\label{fig:traj}
\end{figure}
